\section{Manifiesto de evidencias y declaraci\'on de uso de IA}

\subsection{Manifiesto del audio del ensayo}

La evidencia del roleplay se identifica con el c\'odigo de sesi\'on \texttt{GA-TL-02-ENS-01}
y corresponde al perfil \textbf{SH-01 (Cajero / Vendedor de Mostrador)}. Sus datos se
registran en el repositorio sin publicar el archivo de audio.

\noindent\begin{tabularx}{\linewidth}{|p{0.32\linewidth}|Y|}
\hline
\thdr{Campo} & \thdr{Dato verificado}\\
\hline
Nombre del archivo & \texttt{ensayo\_entrevista\_SH01\_20261120.mp3}\\
\hline
\rowcolor{rowgray} Fecha de grabaci\'on & 20 de noviembre de 2026\\
\hline
Duraci\'on & 10:35\\
\hline
\rowcolor{rowgray} Formato & MP3\\
\hline
SHA-256 & \texttt{\footnotesize 31986071108A8BD445B3DCF5C045BB5EF207288A36967A2E2DDD1CAF397FDDA1}\\
\hline
\end{tabularx}

\vspace{4pt}
La huella se obtuvo con \texttt{Get-FileHash \ldots -Algorithm SHA256} sobre el archivo
final, sin conversiones posteriores. El audio se inicia con el guion de apertura y la
autorizaci\'on verbal del participante, seg\'un el protocolo de consentimiento descrito en
la secci\'on de contexto y en el registro \texttt{ACTA-CONS-01}.

\subsection{Custodia de la evidencia}

De acuerdo con la LOPDP del Ecuador \cite{lopdp2021}, el archivo de audio se conserva en
una carpeta privada de OneDrive institucional (\texttt{@uteq.edu.ec}) del integrante
responsable, compartida con el docente \texttt{gguerrero@uteq.edu.ec}. No se publica en el
repositorio ning\'un dato personal de terceros: ni c\'edulas, tel\'efonos, correos personales,
audio ni consentimientos firmados. El registro detallado se encuentra en
\texttt{GA/TL-02/evidencias/MANIFIESTO.md}.

\subsection{Declaraci\'on de uso de IA como apoyo}

\begin{itemize}
    \item Herramienta: OpenCode, modelo \texttt{Big Pickle} (proveedor \texttt{opencode}).
    \item Versi\'on: modelo \texttt{Big Pickle}, sesi\'on de los d\'ias 20 y 21 de noviembre de 2026.
    \item Prop\'osito: apoyo en la estructura \LaTeX{}, en la revisi\'on de las normas de entrega
          y en el c\'alculo de la huella SHA-256 del audio del ensayo.
    \item Verificaci\'on: el equipo revis\'o manualmente cada secci\'on, contrast\'o los datos del
          manifiesto con el archivo original mediante \texttt{Get-FileHash} y valid\'o la
          compilaci\'on local con MiKTeX. Las respuestas y apreciaciones del roleplay provienen
          de la grabaci\'on real del equipo.
\end{itemize}
